\apendice{Instrumento de Avaliação de Aceitação Tecnológica (TAM)}
\label{ap:instrumento-tam}

Este apêndice reproduz, na íntegra e na ordem original, o questionário de avaliação de aceitação tecnológica aplicado aos participantes da validação descrita no Capítulo~\ref{cap:validacao-tam}. O instrumento foi implementado na plataforma Google Forms sob o título \textit{``Avaliação de Aceitação Tecnológica (TAM) dos Objetivos de Desenvolvimento Sustentável (ODS) em Mossoró-RN''} e disponibilizado aos respondentes entre fevereiro e março de 2026.

\section*{Texto de Apresentação do Formulário}

\begin{citacao}
Prezados(as) Professores(as), Pesquisadores(as) e membros da comunidade acadêmica,

Me chamo João Marinho Neto, sou mestrando em Ciência da Computação na UFERSA, e convido vocês a participarem da validação técnica da minha pesquisa sobre os Objetivos de Desenvolvimento Sustentável (ODS) em Mossoró.

Preciso da sua expertise acadêmica e visão multidisciplinar para avaliar os diagnósticos gerados e verificar a robustez dos resultados apresentados.

\textbf{O que foi feito:} O estudo avalia o desempenho de Mossoró em relação aos ODS 9 (Indústria, Inovação e Infraestrutura) e ODS 16 (Paz, Justiça e Instituições Eficazes). A metodologia utilizou extração de dados públicos (Google Places API) e aplicação de Inteligência Artificial -- especificamente Processamento de Linguagem Natural (PLN) e Análise de Sentimentos -- para processar comentários e mapear desafios urbanos e institucionais. Os resultados finais foram consolidados em painéis interativos (dashboards).

\textbf{Por que sua avaliação importa:} A validação técnica por especialistas de diferentes áreas é uma etapa crucial da pesquisa para: confirmar a precisão e a coerência analítica dos diagnósticos extraídos pelo modelo; avaliar a representatividade dos problemas levantados sob uma ótica científica; garantir o rigor metodológico e a aplicabilidade prática das conclusões geradas.

\textbf{Como participar:} 1. Antes de responder, acesse e explore os painéis interativos neste link: \url{https://dashboard-ods-mossoro.netlify.app/}; 2. Em seguida, preencha o questionário abaixo. O tempo estimado é de 3 a 5 minutos.

Todas as respostas são anônimas, confidenciais e os dados serão utilizados exclusivamente para fins de pesquisa científica.
\end{citacao}

\section*{Bloco 1 -- Identificação}

\begin{enumerate}
    \item E-mail. \textit{(resposta curta, obrigatória)}
    \item Nome do respondente. \textit{(resposta curta, obrigatória)}
    \item Departamento onde atua. \textit{(resposta curta)}
    \item Data da sua avaliação. \textit{(data)}
\end{enumerate}

\section*{Bloco 2 -- Percepção de Utilidade (PU)}

\noindent Itens respondidos em escala Likert de cinco pontos, com as opções: \textit{Discordo Totalmente; Discordo; Neutro; Concordo; Concordo Totalmente} (marcar apenas uma opção).

\begin{enumerate}
    \setcounter{enumi}{4}
    \item \textbf{(PU1)} Percepção de Utilidade (PU): Acredito que utilizar estes dashboards melhorará meu desempenho na tomada de decisão.
    \item \textbf{(PU2)} Percepção de Utilidade (PU): As informações fornecidas pelos dashboards são relevantes para minhas responsabilidades.
    \item \textbf{(PU3)} Percepção de Utilidade (PU): A utilização dos dashboards me permite acessar dados e indicadores de forma mais eficaz.
\end{enumerate}

\section*{Bloco 3 -- Percepção de Facilidade de Uso (PFU)}

\begin{enumerate}
    \setcounter{enumi}{7}
    \item \textbf{(PFU1)} Percepção de Facilidade de Uso (PFU): Acho o uso dos dashboards fácil de aprender.
    \item \textbf{(PFU2)} Percepção de Facilidade de Uso (PFU): Minha interação com os dashboards é clara e compreensível.
    \item \textbf{(PFU3)} Percepção de Facilidade de Uso (PFU): Acho simples navegar e encontrar as informações que preciso nos dashboards.
\end{enumerate}

\section*{Bloco 4 -- Intenção Comportamental de Uso (ICU)}

\begin{enumerate}
    \setcounter{enumi}{10}
    \item \textbf{(ICU1)} Intenção Comportamental de Uso (ICU): Pretendo usar os dashboards regularmente para acompanhar os indicadores da minha área.
    \item \textbf{(ICU2)} Intenção Comportamental de Uso (ICU): Em geral, pretendo usar os dashboards para obter diagnósticos e análises sobre os Objetivos de Desenvolvimento Sustentável (ODS) em Mossoró.
\end{enumerate}

\section*{Bloco 5 -- Avaliação Complementar dos Dashboards}

\begin{enumerate}
    \setcounter{enumi}{12}
    \item Qual dos dois dashboards você considera mais útil para suas atividades diárias? \textit{(lista suspensa)}
    \begin{itemize}
        \item Dashboard Geral (Visão Abrangente dos ODS)
        \item Dashboard Específico (Foco em detalhamento/diagnóstico)
        \item Ambos são igualmente úteis
        \item Nenhum é útil
    \end{itemize}
    \item Em uma escala de 1 a 5, quão satisfeito você está com a clareza e precisão dos dados apresentados nos Dashboards Principais? \textit{(escala de 1 = Muito Insatisfeito a 5 = Muito Satisfeito)}
    \item Em uma escala de 1 a 5, quão satisfeito você está com a clareza e profundidade dos diagnósticos oferecidos na Tela de Diagnósticos? \textit{(escala de 1 = Muito Insatisfeito a 5 = Muito Satisfeito)}
    \item Avalie os seguintes aspectos dos dashboards: \textit{(matriz de avaliação; marcar uma opção por linha, entre Ruim, Regular, Bom e Excelente)}
    \begin{itemize}
        \item Velocidade de carregamento
        \item Design e estética visual
        \item Organização e fluxo de navegação
        \item Capacidade de filtrar e detalhar dados (\textit{drill-down})
    \end{itemize}
    \item Quais informações/indicadores você acredita que estão faltando ou deveriam ser adicionados aos dashboards para melhorar sua utilidade? \textit{(resposta aberta)}
    \item Quais dificuldades ou problemas você encontrou ao utilizar os dashboards (ex: dados incorretos, links quebrados, lentidão, confusão)? \textit{(resposta aberta)}
    \item Marque as opções que melhor descrevem a utilidade dos diagnósticos gerados pela IA fornecidos pelos dashboards: \textit{(múltipla seleção)}
    \begin{itemize}
        \item Ajudam a identificar áreas críticas para intervenção
        \item São fundamentais para o planejamento estratégico da minha área
        \item Oferecem uma visão clara do progresso de Mossoró nos Objetivos de Desenvolvimento Sustentável (ODS)
        \item São muito superficiais e pouco acionáveis
        \item Não utilizo os diagnósticos
    \end{itemize}
    \item Se tivéssemos que ranquear a importância (do mais importante para o menos importante) dos dashboards para sua função, qual seria sua ordem? \textit{(matriz de avaliação; marcar uma opção por linha, entre Mais Importante e Menos Importante)}
    \begin{itemize}
        \item Informações dos Dashboards Principais
        \item Detalhes e Diagnósticos Gerados pela IA
    \end{itemize}
    \item Comentários finais: \textit{(resposta aberta)}
\end{enumerate}
